\chapter{Mathematical background}
\label{chap:background}

This chapter assumes familiarity with the finite element idea: approximate a
function in a finite-dimensional space and require a residual to vanish against
test functions. DG retains that idea but removes continuity between elements.
The price is that communication across each interface must be defined explicitly.
We first explain this communication, then the time update that the filters modify. The introductory DG treatment follows the finite element viewpoint of~\cite{HesthavenWarburton2008}.

\section{Conservation laws and the quantity being conserved}
\label{sec:rkdg}
A conservation law describes a balance between local storage and transport:
\begin{equation}
 \partial_t\boldsymbol u+\nabla\cdot\boldsymbol f(\boldsymbol u)=0,
 \qquad \boldsymbol u(\boldsymbol x,0)=\boldsymbol u_0(\boldsymbol x).
 \label{eq:conservation-law}
\end{equation}
Here $\boldsymbol u\in\R^m$ contains the conserved quantities, $\Omega$ is the
spatial domain, and $\boldsymbol f$ gives their fluxes in the $d$ coordinate
directions. Boundary conditions supply the missing information on the exterior
of $\Omega$. Integrating over a region says that its total content changes only
through flux across the region's boundary. We discuss the scalar case first;
the same balance applies component by component to a system.

\section{What changes from continuous FEM?}
Let $\mathcal T_h$ be a mesh of elements $K$. In a continuous finite element
space, neighbouring polynomial pieces must agree on a shared interface. In DG,
each element has its own coefficients, so that agreement is not imposed. At a
face, $u_h^-$ denotes the value approached from inside the current element and
$u_h^+$ the value approached from the neighbour. These are \emph{traces}, not
positive and negative parts of a number. Their difference is a \emph{jump}.

\begin{figure}[htbp]
\centering
\begin{tikzpicture}[>=Latex,scale=1.0]
 \draw[thick] (0,0) rectangle (3,1.6);
 \draw[thick] (3,0) rectangle (6,1.6);
 \node at (1.5,1.25) {element $K$};\node at (4.5,1.25) {neighbour $N$};
 \draw[blue!70!black,thick] (.2,.4)--(2.8,.95);
 \draw[orange!80!black,thick] (3.2,.4)--(5.8,.65);
 \node[below] at (2.55,.8) {$u_h^-$};\node[above] at (3.6,.45) {$u_h^+$};
 \draw[->,thick] (3,1.8)--(4,1.8);\node[above] at (3.5,1.8) {$n_K$};
 \node[below,align=center] at (3,-.25) {independent traces; one shared numerical flux};
\end{tikzpicture}
\caption{A schematic pair of elements. The separate polynomial pieces may
jump. A numerical flux gives both elements a consistent exchange across their
common face. The drawn heights illustrate values, not element geometry.}
\label{fig:dg-traces}
\end{figure}

This locality makes DG convenient for discontinuous solutions and distributed
meshes. It does not by itself prevent oscillations: a high-order polynomial
can overshoot near a shock even when its cell average is reasonable.

\subsection{Spaces, coefficients, and geometry}
Each physical element is the image of a reference element $\widehat K$ under
an invertible map $T_K$. Its order-$k$ space is
\begin{equation}
 V^k(K)=\{\widehat v\circ T_K^{-1}:\widehat v\in\widehat V^k(\widehat K)\},
 \qquad V_h^k=\{v:v|_K\in V^k(K),\ K\in\mathcal T_h\}.
 \label{eq:ofdg-space}
\end{equation}
The reference family is total-degree $P^k$ on a simplex, tensor-product $Q^k$
on a quadrilateral or hexahedron, and the corresponding triangle--segment
product on a prism. Thus ``order $k$'' does not always mean total degree $k$.
A curved map generally produces functions that are not polynomials in physical
coordinates. All lower-order spaces in this work use the same map, so they
remain nested: $V^0(K)\subset\cdots\subset V^k(K)$.

With basis functions $\phi_i$, the approximation is
\begin{equation}
 u_h|_K=\sum_{i=1}^{N_K}U_{K,i}\phi_i.
 \label{eq:ofdg-local-expansion}
\end{equation}
$N_K$ is the number of local coefficients, or degrees of freedom (DOFs).
Coefficients may be nodal values or modal expansion coefficients. The function
represented by them, rather than the storage basis, is the object to which
our projections and filters apply.

\section{One element's weak form}
Multiply the scalar conservation law by a test function $w_h\in V^k(K)$,
integrate over $K$, and integrate the spatial divergence by parts:
\begin{equation}
 \int_K\partial_tu_h w_h\,dx
 -\int_K\boldsymbol f(u_h)\cdot\nabla w_h\,dx
 +\int_{\partial K}\widehat f(u_h^-,u_h^+,\boldsymbol n_K)w_h^-\,ds=0.
 \label{eq:rkdg-weak-form}
\end{equation}
The first term is storage; the second accounts for variation inside the element;
the third is exchange across its boundary. $\boldsymbol n_K$ is the outward
unit normal. Since there are two traces, evaluating the physical flux from
only one trace would leave the exchange ambiguous. The \emph{numerical flux}
$\widehat f$ resolves that ambiguity.

The shared baseline uses the Rusanov flux
\begin{equation}
 \widehat f(u^-,u^+,\boldsymbol n)
 =\tfrac12\bigl(\boldsymbol f(u^-)+\boldsymbol f(u^+)\bigr)\cdot\boldsymbol n
   -\tfrac12\lambda(u^+-u^-),
 \label{eq:rusanov}
\end{equation}
where $\lambda$ bounds the relevant normal propagation speeds. The second term
adds interface dissipation proportional to the jump. Equal traces recover the
physical flux. Swapping traces and reversing the normal reverses the numerical
flux, so an internal interface adds to one element exactly what it removes
from the other, subject to consistent quadrature.

Choosing $w_h=1$ gives the cell balance
\begin{equation}
 \frac{d}{dt}\int_K u_h\,dx
 =-\int_{\partial K}\widehat f\,ds.
 \label{eq:cell-balance}
\end{equation}
Internal contributions cancel in the global sum. Domain totals therefore stay
constant for periodic problems, but can change through open boundaries.
A conservative filter must preserve the cell integral established by this DG
update; it does not remove the boundary flux from this balance.

\subsection{Example: linear advection}
For $u_t+a u_x=0$ with constant $a>0$, the physical flux is $au$. Setting
$\lambda=|a|$ in \cref{eq:rusanov} gives $\widehat f=a u^-$ at the right face:
the upstream trace supplies the transported quantity. If two traces differ,
the method communicates their effect through this flux rather than making the
polynomials continuous. This example illustrates the general construction;
it is not an additional governing assumption for OFDG.

\section{Time stepping and the location of filtering}
Substitution of the basis expansion into the weak form gives a finite system
\begin{equation}
 M\frac{dU}{dt}=R(U),
 \qquad M_{ij}=\int_K\phi_i\phi_j\,dx\quad\hbox{on each element}.
\end{equation}
The mass matrix $M$ is block-local; the residual $R$ includes neighbour traces.
An explicit Runge--Kutta method evaluates this residual at several intermediate
\emph{stages}. The implementation provides classical fourth-order RK (RK4), with a CFL
step-size rule depending on mesh scale, order, and characteristic speed.
CFL refers to the Courant--Friedrichs--Lewy restriction: the timestep limits
how far information travels relative to a cell width, with a stricter
restriction at higher approximation order. A common CFL number therefore does
not necessarily produce the same timestep for different evolving states.
RK4's formal ODE order is not a proof of fourth-order accuracy for a nonlinear
filtered shock calculation.

The terminology RKDG means DG in space combined with Runge--Kutta time stepping
\cite{LuEtAl2021,LiuEtAl2022}. In this work, adapted OFDG is normally
applied after a complete step, whereas the OEDG comparator filters each RK
stage. Their cadence is part of the comparison. Exact solution of the damping
substep, introduced next, does not make the combined DG/filter evolution exact.

\section{Nonlinear scalar equations and Euler}
Burgers' equation uses $f(u)=u^2/2$, so its propagation speed depends on the
state; initially smooth data can develop a shock. Compressible Euler has
$m=d+2$ conserved components,
\begin{equation}
 \boldsymbol u=(\rho,\rho\boldsymbol v,E)^\mathsf T,
 \qquad p=(\gamma-1)\left(E-\tfrac12\rho|\boldsymbol v|^2\right).
 \label{eq:euler-state}
\end{equation}
Here $\rho$ is density, $\boldsymbol v$ velocity, $E$ total energy per unit
volume, $p$ pressure, and $\gamma=1.4$ the ratio of specific heats used by the Euler implementation.
The flux in coordinate direction $j$ is
\begin{equation}
 \boldsymbol f_j=(\rho v_j,\ \rho v_j\boldsymbol v+p\boldsymbol e_j,\ (E+p)v_j)^\mathsf T,
 \label{eq:euler-flux}
\end{equation}
with $\boldsymbol e_j$ the coordinate unit vector. The maximum normal wave
speed is $|\boldsymbol v\cdot\boldsymbol n|+c$, where
$c=\sqrt{\gamma p/\rho}$ is the sound speed. This differs from the signed
fluid transport velocity used by the KXRCF inflow decision. Density and pressure
must remain positive for these states and speeds to be physically admissible.
